\documentclass[11pt,a4paper]{article}
\usepackage[T1]{fontenc}
\usepackage[utf8]{inputenc}
\usepackage{lmodern,amsmath,amssymb}
\usepackage[margin=25mm]{geometry}
\usepackage[hidelinks]{hyperref}
\hypersetup{pdftitle={When an action plateau controls a spectral gap},pdfauthor={navstokgap research working notes}}
\newcommand{\tightlist}{\setlength{\itemsep}{0pt}\setlength{\parskip}{0pt}}
\setlength{\emergencystretch}{3em}
\setcounter{secnumdepth}{0}
\begin{document}
\hypertarget{when-an-action-plateau-controls-a-spectral-gap}{%
\section{When an action plateau controls a spectral
gap}\label{when-an-action-plateau-controls-a-spectral-gap}}

A complete collection of observables with bounded total susceptibility
gives a lower relaxation-gap bound. A single positive velocity plateau
instead gives an upper bound, and can miss an arbitrarily slow internal
mode. We prove both statements for finite reversible dynamics and
specify the energy normalization needed for the Dirac comparison.

\hypertarget{operator-observable-and-units}{%
\subsection{1. Operator, observable and
units}\label{operator-observable-and-units}}

Let \(Q\) generate an irreducible continuous-time Markov chain on a
finite set of \(n\ge2\) states, with stationary probabilities
\(\pi_i>0\) and detailed balance \(\pi_iQ_{ij}=\pi_jQ_{ji}\). The
generator acts on functions, with row sums zero. On \(L^2(\pi)\) use
\(\langle f,g\rangle_\pi=\sum_i\pi_i\overline{f_i}g_i\). The centered
space \(\mathcal H_0=\{f:\langle1,f\rangle_\pi=0\}\) has dimension
\(n-1\). All domains are the entire indicated finite-dimensional space.

Set \(A=-Q|_{\mathcal H_0}\). It is positive self-adjoint with
eigenvalues \(0<\gamma_1\le\cdots\le\gamma_{n-1}\). The full relaxation
gap is \(\gamma_1\), in inverse-time units. For a real centered
observable \(v\) with velocity units, define

\[
\chi(v)=\langle v,A^{-1}v\rangle_\pi
=\int_0^\infty\langle v,e^{tQ}v\rangle_\pi\,dt,
\qquad H(v)=2m\chi(v),\quad m>0.
\]

\(\chi\) has units length squared per time, and \(H\) has action units.
This is C019's long-observation plateau for the integrated velocity. The
Green--Kubo/Poisson representation is matched to Pavliotis in
\href{../references/batches/B20.md}{B20}; the finite-dimensional
consequences below are derived explicitly. Time here is the chain's
evolution parameter.

\hypertarget{one-observable-the-product-and-its-direction-c041}{%
\subsection{2. One observable: the product and its direction
(C041)}\label{one-observable-the-product-and-its-direction-c041}}

Write \(v=\sum_j a_je_j\) in an orthonormal eigenbasis of \(A\), and
\(\sigma_v^2=\|v\|_\pi^2>0\). Then

\[
H(v)=2m\sum_j\frac{|a_j|^2}{\gamma_j},\qquad
H(v)\gamma_1\le2m\sigma_v^2\le H(v)\gamma_{n-1}.
\]

\textbf{Proof.} Integrate each exponential in
\(\langle v,e^{tQ}v\rangle_\pi=\sum_j|a_j|^2e^{-\gamma_jt}\) and use
\(\sum_j|a_j|^2=\sigma_v^2\). The left inequality gives
\(\gamma_1\le2m\sigma_v^2/H(v)\): an upper gap bound. The normalized
integral time \(\tau_v=\chi(v)/\sigma_v^2\) is a weighted average of
\(1/\gamma_j\); in particular \(\tau_v\le1/\gamma_1\).

For the two-state chain \(Q=\lambda(\sigma_x-I)\) and \(v=u(1,-1)\),
\(\pi=(1/2,1/2)\), with \(u,\lambda>0\), the centered space is
one-dimensional. Therefore

\[
\gamma_1=2\lambda,\qquad H(v)=\frac{mu^2}{\lambda},\qquad
H(v)\gamma_1=2mu^2.
\]

Here \(\sigma_x\) exchanges the two signs. This equality identifies the
source of the two-state product: every centered mode is visible to
velocity.

\hypertarget{fixed-plateau-with-a-closing-full-gap-c041}{%
\subsection{3. Fixed plateau with a closing full gap
(C041)}\label{fixed-plateau-with-a-closing-full-gap-c041}}

Take two independent signs \(s,r\in\{-1,1\}\) with rates \(\lambda>0\)
and \(\epsilon>0\), respectively, and define velocity \(v(s,r)=us\),
with fixed \(0<u<c\). The four-state generator is

\[
Q_{\lambda,\epsilon}=\lambda(\sigma_x-I)\otimes I
+I\otimes\epsilon(\sigma_x-I).
\]

The uniform measure is stationary and reversible. The functions
\(1,s,r,sr\) form an orthonormal eigenbasis with eigenvalues of \(-Q\)
equal to \(0,2\lambda,2\epsilon,2(\lambda+\epsilon)\). Velocity overlaps
only the \(s\) mode. Thus

\[
H(v)=\frac{mu^2}{\lambda},\qquad
\gamma_1=2\min(\lambda,\epsilon)\longrightarrow0
\quad\hbox{as }\epsilon\downarrow0.
\]

Every positive-\(\epsilon\) member is irreducible, with the same bounded
speed and the same positive plateau. At the limiting parameter the label
becomes conserved. The velocity path law itself is unchanged throughout
the family. The missing information is the label's relaxation:
finite-rate assumptions for each member give positivity member by
member, while a uniform lower gap requires additional control over the
family.

\hypertarget{a-sufficient-condition-complete-observability-c042}{%
\subsection{4. A sufficient condition: complete observability
(C042)}\label{a-sufficient-condition-complete-observability-c042}}

Choose centered velocity-unit observables \(f_1,\ldots,f_r\) and suppose
that, for every \(g\in\mathcal H_0\),

\[
\sum_{a=1}^r|\langle f_a,g\rangle_\pi|^2
\ge\alpha\|g\|_\pi^2,\qquad\alpha>0.
\]

This \emph{frame bound} means that the collection detects every centered
mode; \(\alpha\) has velocity-squared units. Put \(S=\sum_a\chi(f_a)\).
Then

\[
\boxed{\gamma_1\ge\frac{\alpha}{S}.}
\]

If \(H(f_a)\le B_a\), it follows that \(\gamma_1\ge2m\alpha/\sum_aB_a\).

\textbf{Proof.} Set \(b_j=\sum_a|\langle e_j,f_a\rangle_\pi|^2\). The
frame hypothesis gives \(b_j\ge\alpha\) for every unit eigenvector.
Hence

\[
S=\sum_j\frac{b_j}{\gamma_j}\ge\frac{\alpha}{\gamma_1}.
\]

Multiplication by \(\gamma_1/S\) proves the result. Equivalently, the
frame operator \(F=\sum_a|f_a\rangle\langle f_a|\) satisfies
\(F\ge\alpha I\) and \(S=\operatorname{Tr}(A^{-1}F)\).

For a family, uniform \(\alpha\ge\alpha_0>0\) and \(S\le S_0<\infty\)
give the uniform bound \(\gamma_1\ge\alpha_0/S_0\). When masses vary,
the action-valued version must control \(\sum_aB_a/(2m)\), rather than
merely \(\sum_aB_a\). The frame can have full rank only if \(r\ge n-1\).
These requirements expose the cost of carrying the estimate to larger
systems.

The exact all-observable formulation is

\[
\sup_{f\in\mathcal H_0\setminus\{0\}}
\frac{\chi(f)}{\|f\|_\pi^2}=\|A^{-1}\|=\frac1{\gamma_1}.
\]

The spectral expansion proves the upper bound, and a slowest eigenvector
attains equality. It is the inverse-operator version of a Poincaré
bound.

In the four-state example choose \(f_1=us\), \(f_2=ur\), \(f_3=usr\).
They give \(F=u^2I\) on \(\mathcal H_0\), and

\[
S=u^2\left[\frac1{2\lambda}+\frac1{2\epsilon}
+\frac1{2(\lambda+\epsilon)}\right].
\]

The complete collection detects the approaching gap closure through
\(\chi(f_2)\to\infty\). Each observable stays bounded in magnitude by
\(u\). Thus speed control and susceptibility control address different
premises.

\hypertarget{energy-normalization-and-the-companion-field-problem}{%
\subsection{5. Energy normalization and the companion field
problem}\label{energy-normalization-and-the-companion-field-problem}}

On full \(L^2(\pi)\) define \(\mathcal E=K(-Q)\) with a supplied action
unit \(K>0\). Constants form its unique zero-energy ground space; its
excitation gap is \(\Delta_{\mathcal E}=K\gamma_1\). In the two-state
model, setting \(K=H(v)\) gives

\[
\Delta_{\mathcal E}=2mu^2.
\]

The numerical substitution \(u=c\) matches C039's Dirac rest-branch
separation \(2mc^2\). The two operators have different state spaces and
spectral meanings: \(\mathcal E\) is a nonnegative finite-state
relaxation operator, while the Dirac operator has positive and negative
single-particle branches. A quantum field vacuum gap additionally
requires its physical Hilbert space and the continuum/infinite-volume
construction.

For the companion programme, C042 supplies a candidate form of estimate:
observable coverage bounded below, inverse-generator response bounded
above. Transferring it requires identifying the physical generator and
proving that both bounds survive the relevant limits. An algorithmic
sampling chain's relaxation time is a distinct object from physical
Euclidean time evolution.

\hypertarget{next-test-and-reproduction}{%
\subsection{6. Next test and
reproduction}\label{next-test-and-reproduction}}

G02 asks whether mechanical composition preserves a useful observable
collection and its quantitative coverage. Start with the four-state
example, then weaken access to the label. Track the smallest frame
eigenvalue and total susceptibility together; a gap bound based on
either one alone loses the information isolated above. A09 continues the
independent action-scale and coherent-composition selection task.

Run \texttt{python3\ scripts/susceptibility\_gap\_checks.py}. The
\href{../docs/batches/B20/susceptibility-source-companion.md}{B20
companion} records the Green--Kubo source match and the bounded
prior-art coverage. C041 is a spectral/product-chain consequence of
C019; C042 is an elementary finite-dimensional observability criterion,
with no novelty claim.
\end{document}
